% SPDX-License-Identifier: LPPL-1.3c
% Copyright (C) 2026 Christophe Jorssen
% Author and Current Maintainer: Christophe Jorssen <christophe.jorssen@gmail.com>
% LPPL maintenance status: maintained
% This file is part of luafemm. See LICENSE and LICENSES.md for its terms.
\documentclass[varwidth=150mm,border=3mm]{standalone}
\usepackage{tikz}
\usetikzlibrary{femm}
\begin{document}
\tikzset{femm/materials/gap iron/.style={mu r=100}}
\begin{tikzpicture}[
  femm/problem={xmin=-8,xmax=8,ymin=-8,ymax=8,mesh size=1.2,depth=10},
  femm/sliding band={name=native-ring,method=femm5,
    inner radius=3,outer radius=3.6,segments=64},
  scale=2.5,baseline=(current bounding box.center)]
  \draw[femm/region={material=gap iron}]
    (0,0) ellipse[x radius=2.6,y radius=1.2];
  \draw[femm/region={material=copper,current density=1e7}]
    (-6,-1) rectangle (-5,1);
  \draw[femm/region={material=copper,current density=-1e7}]
    (5,-1) rectangle (6,1);
\end{tikzpicture}
\begin{tikzpicture}[
  femm/problem={sliding band=native-ring,rotor angle=25},
  scale=2.5,baseline=(current bounding box.center)]
  \pic {femm geometry};
  \pic[femm/lines=13] {femm field};
  \femmexport[file=\jobname-native-gap.fem,format=fem]
  \femmexport[file=\jobname-native-gap.ans,format=ans]
\end{tikzpicture}

\par
\begin{tikzpicture}[
  femm/import={file=\jobname-native-gap.ans,
    age interpolation=femm5},scale=2.5]
  \pic {femm geometry};
  \pic[femm/lines=13] {femm field};
\end{tikzpicture}
\end{document}
